% Part: normal-modal-logic
% Chapter: tableaux
% Section: introduction

\documentclass[../../../include/open-logic-section]{subfiles}

\begin{document}

\olfileid{nml}{tab}{int}

\olsection{Pambuka}

!!^{tableau}s iku wit tartamtu kang cabange tumuju mangisor lan dumadi
saka !!{signed
  formula}s, yaiku pasangan saka tandha biji-kabeneran
($\True$ utawa $\False$) lan !!a{sentence}
\[
\sFmla{\True}{!A} \text{ utawa } \sFmla{\False}{!A}.
\]
!!^a{tableau} diwiwiti kanthi sawetara \emph{asumsi}. Saben
!!{signed formula} sabanjure diasilake kanthi ngetrapake salah siji
aturan inferensi. Sawetara aturan inferensi nambah siji utawa luwih
!!{signed formula}s ing pucuk wit; aturan liyane nambah rong pucuk
anyar, saengga ngasilake rong cabang. Aturan-aturan mau ngasilake
!!{signed formula}s kang !!{formula}-ne luwih prasaja tinimbang
kang ana ing !!{signed formula} kang dienggo ngetrapake aturan iku.
Yen sawijining cabang ngemot $\sFmla{\True}{!A}$ lan
$\sFmla{\False}{!A}$, cabang iku diarani \emph{katutup}. Yen saben
cabang ing !!a{tableau} katutup, kabeh !!{tableau} iku katutup.
!!{tableau} katutup mujudake !!a{derivation} kang nuduhake yen
himpunan !!{signed formula}s kang digunakake kanggo miwiti
!!{tableau} iku ora bisa dicukupi. Iki bisa digunakake kanggo
netepake relasi~$\Proves$:
$\Gamma \Proves !A$ yen lan mung yen ana himpunan winates~$\Gamma_0 = \{!B_1,
\dots, !B_n\} \subseteq \Gamma$ kang nduweni !!{tableau} katutup
kanggo asumsi
\[
\{\sFmla{\False}{!A}, \sFmla{\True}{!B_1}, \dots, \sFmla{\True}{!B_n}\}.
\]

Kanggo logika modal, kita kudu ngambakake pangerten !!{signed
formula} lan nambah aturan kang nyakup~\iftag{prvBox}{$\Box$\iftag{prvDiamond}{ lan
    $\Diamond$}}{$\Diamond$}. Saliyane tandha ($\True$ utawa
$\False$), !!{formula}s ing !!{tableau}s modal uga nduweni
\emph{prefiks}~$\sigma$. Prefiks iku urutan bilangan bulat positif
kang ora kosong, yaiku $\sigma \in (\PosInt)^* \setminus
\{\emptyseq\}$. Kita nulis prefiks kasebut tanpa kurung
$\tuple{\ }$ lan misahake saben !!{element}s nganggo~$.$,
dudu $,$. Yen $\sigma$ sawijining prefiks, $\sigma.n$ padha karo $\sigma
\concat \tuple{n}$; upamane, yen $\sigma = 1.2.1$, $\sigma.3$ iku
$1.2.1.3$. Dadi, upamane,
\[
\sFmla{\True}{\Box !A \lif !A}[1.2]
\]
iku \emph{!!{signed formula} mawa prefiks} (utawa cekake
\emph{!!{formula} mawa prefiks}).

Miturut intuisi, prefiks iku menehi jeneng marang jagad ing sawijining
model kang bisa nyukupi !!{formula}s ing sawijining cabang saka
!!a{tableau}; yen $\sigma$ menehi jeneng marang sawijining jagad,
mula $\sigma.n$ menehi jeneng marang jagad kang bisa diakses saka
jagad kang dijenengi~$\sigma$.

\end{document}
